\documentclass[12pt]{article}
\usepackage[utf8]{inputenc}
\usepackage[T1]{fontenc}
\usepackage[margin=1in]{geometry}
\usepackage{setspace}
\usepackage{xcolor}
\usepackage[colorlinks=true, urlcolor=blue!60!black, linkcolor=blue!60!black]{hyperref}
\setlength{\parindent}{2em}
\setlength{\parskip}{1.25em}
\renewcommand{\baselinestretch}{1.0}
\begin{document}
\begin{center}
{\LARGE \textit{Regents of the University of California v. Bakke}, 438 U.S. 265 (1978)\par}
\vspace{0.8em}
{\large Case Brief\par}
\end{center}
\vspace{1em}

\section{Memory Jogger}
In a deeply fractured 4--1--4 decision anchored by Justice Powell's pivotal opinion, the Supreme Court struck down UC Davis Medical School's rigid 16-seat racial quota and ordered Allan Bakke admitted, while simultaneously holding that public universities could constitutionally consider race as one ``plus'' factor among many to achieve student-body diversity (\href{https://www.courtlistener.com/opinion/109930/regents-of-the-university-of-california-v-bakke/}{CourtListener}; \href{https://supreme.justia.com/cases/federal/us/438/265/}{Justia}).

\section{Detailed Case Facts}
The Medical School of the University of California at Davis opened in 1968 with an entering class of 50 students and expanded to 100 students in 1971 (\href{https://supreme.justia.com/cases/federal/us/438/265/}{438 U.S. at 272}). Because its first class included three Asian students but no Black, Mexican-American, or American Indian students, the faculty created a two-track admissions system consisting of a regular admissions program and a special admissions program (\href{https://supreme.justia.com/cases/federal/us/438/265/}{438 U.S. at 272--273}).

Under the \textbf{regular admissions program}, applicants with an undergraduate GPA below 2.5 on a 4.0 scale were summarily rejected (\href{https://supreme.justia.com/cases/federal/us/438/265/}{438 U.S. at 273}). Roughly one in six regular applicants received an interview, after which committee members rated each applicant on a 1--100 scale based on GPAs, MCAT scores, interview summaries, recommendations, and extracurricular activities to produce a ``benchmark'' score (out of 500 in 1973 with five raters, and out of 600 in 1974 with six raters) (\href{https://supreme.justia.com/cases/federal/us/438/265/}{438 U.S. at 273--274}).

Under the \textbf{special admissions program}, a separate committee (a majority of whom were members of minority groups) screened applicants who indicated on the 1973 form that they were ``economically and/or educationally disadvantaged'' or on the 1974 form that they belonged to a ``minority group'' (Blacks, Chicanos, Asians, or American Indians) (\href{https://supreme.justia.com/cases/federal/us/438/265/}{438 U.S. at 274}). Special applicants were exempt from the 2.5 GPA cutoff and were never rated or compared against regular applicants; the special committee recommended candidates until 16 of the 100 seats were filled (\href{https://supreme.justia.com/cases/federal/us/438/265/}{438 U.S. at 275}). From 1971 through 1974, 63 minority students were admitted through the special program (21 Black, 30 Mexican-American, 12 Asian) and 44 through the regular program (1 Black, 6 Mexican-American, 37 Asian); although 73 disadvantaged white students applied to the special program in 1973 and 172 in 1974, none was admitted (\href{https://supreme.justia.com/cases/federal/us/438/265/}{438 U.S. at 275--276 \& n.5}).

Allan Bakke, a white male engineer in his thirties, applied to UC Davis Medical School in 1973 and 1974 under the regular program (\href{https://supreme.justia.com/cases/federal/us/438/265/}{438 U.S. at 276}; \href{https://www.oyez.org/cases/1979/76-811}{Oyez}). Bakke had a 3.44 science GPA, a 3.46 overall GPA, and MCAT percentiles of 96 (Verbal), 94 (Quantitative), 97 (Science), and 72 (General Information) (\href{https://supreme.justia.com/cases/federal/us/438/265/}{438 U.S. at 277 n.7}). In 1973, he applied late and received a strong benchmark score of 468/500, but was rejected because no regular applicants below 470 were then being admitted, even though four special admissions seats remained unfilled (\href{https://supreme.justia.com/cases/federal/us/438/265/}{438 U.S. at 276}). After protesting the special program in a letter to Associate Dean Dr. George H. Lowrey, Bakke reapplied early in 1974; he scored 549/600 (including a low rating of 86 from Dr. Lowrey, who interviewed him) and was rejected again without being placed on the waiting list (\href{https://supreme.justia.com/cases/federal/us/438/265/}{438 U.S. at 276--277}). In both years, special admittees entered with substantially lower average GPAs (2.88 overall in 1973; 2.62 overall in 1974) and MCAT percentiles (35th and 37th in Science) than Bakke (\href{https://supreme.justia.com/cases/federal/us/438/265/}{438 U.S. at 277 \& n.7}).

\section{Procedural History}
\begin{enumerate}
\item \textbf{Superior Court of California (Yolo County):} Bakke sued for mandatory, injunctive, and declaratory relief compelling his admission, alleging violations of the Equal Protection Clause of the Fourteenth Amendment, Article I, \S~21 (later Art. I, \S~7) of the California Constitution, and \S~601 of Title VI of the Civil Rights Act of 1964, 42 U.S.C. \S~2000d (\href{https://supreme.justia.com/cases/federal/us/438/265/}{438 U.S. at 277--278 \& nn.9--11}). The University cross-complained for a declaratory judgment upholding the program (\href{https://supreme.justia.com/cases/federal/us/438/265/}{438 U.S. at 278}). The trial court held that the special program operated as an unlawful racial quota violating the Federal Constitution, the California Constitution, and Title VI, and enjoined the University from considering race in admissions, but refused to order Bakke's admission because he had not proved he would have been admitted absent the program (\href{https://supreme.justia.com/cases/federal/us/438/265/}{438 U.S. at 278--279}).
\item \textbf{Supreme Court of California (\textit{Bakke v. Regents of the University of California}, 18 Cal. 3d 34, 553 P.2d 1152 (1976)):} Transferred the case directly from the trial court (\href{https://supreme.justia.com/cases/federal/us/438/265/}{438 U.S. at 279}). Applying strict scrutiny solely under the Fourteenth Amendment (without reaching state constitutional or Title VI grounds), the court ruled 6--1 (Justice Stanley Mosk writing; Justice Matthew O. Tobriner dissenting) that the special program was not the least intrusive means of achieving its goals and affirmed the injunction barring any consideration of race (\href{https://supreme.justia.com/cases/federal/us/438/265/}{438 U.S. at 279--280}; \href{https://en.wikipedia.org/wiki/Regents_of_the_University_of_California_v._Bakke}{Wikipedia}). The court shifted the burden of proof to the University; when the University conceded on rehearing that it could not prove Bakke would have been rejected without the special program, the court amended its judgment to order Bakke's admission (\href{https://supreme.justia.com/cases/federal/us/438/265/}{438 U.S. at 280--281}).
\item \textbf{Supreme Court of the United States:} The Court stayed the admission order (429 U.S. 953 (1976)), granted certiorari (429 U.S. 1090 (1977)), heard argument on October 12, 1977, requested supplementary briefing on Title VI (434 U.S. 900 (1977)), and issued its fractured judgment on June 28, 1978, affirming in part and reversing in part (\href{https://supreme.justia.com/cases/federal/us/438/265/}{438 U.S. at 265, 281}).
\end{enumerate}

\section{Judicial Votes}
\textbf{Vote Breakdown:} Iconic \textbf{4--1--4} split producing \textbf{six separate opinions}, with no single opinion commanding a majority in full (\href{https://supreme.justia.com/cases/federal/us/438/265/}{438 U.S. at 267--268}; \href{https://www.law.cornell.edu/supremecourt/text/438/265}{Cornell LII}).
\begin{itemize}
\item \textbf{Pivotal Swing Opinion (1 Justice):} \textbf{Justice Lewis F. Powell, Jr.} announced the judgment of the Court (\href{https://supreme.justia.com/cases/federal/us/438/265/}{438 U.S. at 269}). Justices Brennan, White, Marshall, and Blackmun joined Parts I and V-C of Powell's opinion (making 5 votes for the factual background and for reversing the California injunction against any use of race), and Justice White also joined Part III-A (\href{https://supreme.justia.com/cases/federal/us/438/265/}{438 U.S. at 267, 272 n.**}). Powell spoke for himself alone in Parts II, III-B, III-C, IV, V-A, V-B, and VI.
\item \textbf{Affirming Bakke's Admission and Striking Down the 16-Seat Quota (5--4):} \textbf{Powell} (on Fourteenth Amendment grounds) joined with the \textbf{Stevens Bloc} (\textbf{Justice John Paul Stevens}, joined by \textbf{Chief Justice Warren E. Burger} and \textbf{Justices Potter Stewart and William H. Rehnquist}, concurring in the judgment in part and dissenting in part on Title VI statutory grounds) (\href{https://supreme.justia.com/cases/federal/us/438/265/}{438 U.S. at 267--268, 271, 408}).
\item \textbf{Reversing the Blanket Ban on Considering Race and Holding Title VI Coextensive with Equal Protection (5--4):} \textbf{Powell} joined with the \textbf{Brennan Bloc} (\textbf{Justices William J. Brennan, Jr., Byron R. White, Thurgood Marshall, and Harry A. Blackmun}, concurring in the judgment in part and dissenting in part) (\href{https://supreme.justia.com/cases/federal/us/438/265/}{438 U.S. at 267--268, 272, 324}).
\item \textbf{Additional Separate Opinions:} Within the Brennan Bloc, \textbf{Justice White} (438 U.S. at 379), \textbf{Justice Marshall} (438 U.S. at 387), and \textbf{Justice Blackmun} (438 U.S. at 402) each filed separate opinions (\href{https://supreme.justia.com/cases/federal/us/438/265/}{438 U.S. at 268}).
\end{itemize}

\section{Holding}
\textbf{Judgment:} Affirmed in part (invalidating UC Davis's special admissions program and ordering Allan Bakke admitted) and reversed in part (vacating the California Supreme Court's injunction forbidding the University from taking race into account as a factor in future admissions decisions) (\href{https://supreme.justia.com/cases/federal/us/438/265/}{438 U.S. at 266--267, 272}).

\textbf{Rule Announced:} Two distinct five-Justice majorities established that: (1) Title VI of the Civil Rights Act of 1964 proscribes only those racial classifications that would violate the Equal Protection Clause if employed by a State (Powell plus Brennan, White, Marshall, and Blackmun) (\href{https://supreme.justia.com/cases/federal/us/438/265/}{438 U.S. at 267, 287, 328}); (2) a state university's rigid racial quota reserving a fixed number of seats exclusively for minority applicants is unlawful and warrants admission of an excluded plaintiff when the university cannot prove he would have been rejected anyway (Powell plus Stevens, Burger, Stewart, and Rehnquist) (\href{https://supreme.justia.com/cases/federal/us/438/265/}{438 U.S. at 267, 320, 421}); and (3) a state university is not categorically barred from considering race as one factor in future admissions decisions (Powell plus Brennan, White, Marshall, and Blackmun) (\href{https://supreme.justia.com/cases/federal/us/438/265/}{438 U.S. at 267, 320, 326}).

\textbf{Holding vs. Dicta:}
\begin{itemize}
\item \textbf{Holding:} The two-part judgment (striking down UC Davis's 16-seat set-aside while lifting the blanket ban on considering race) and the coextensiveness of Title VI with the Equal Protection Clause commanded five votes (\href{https://supreme.justia.com/cases/federal/us/438/265/}{438 U.S. at 267}).
\item \textbf{Contested Precedential Status (Powell's Diversity Rationale):} No other Justice joined Justice Powell's view in Parts IV-D and V-A that ``student body diversity'' rooted in First Amendment academic freedom is a compelling interest justifying a Harvard-style ``plus'' factor under strict scrutiny (\href{https://supreme.justia.com/cases/federal/us/438/265/}{438 U.S. at 311--319}). For 25 years, lower courts split over whether Powell's solo opinion was binding under \textit{Marks v. United States}, 430 U.S. 188 (1977), until the Supreme Court formally endorsed Powell's diversity rationale in \href{https://en.wikipedia.org/wiki/Grutter_v._Bollinger}{\textit{Grutter v. Bollinger}, 539 U.S. 306, 325 (2003)}, before effectively repudiating race-conscious university admissions in \href{https://en.wikipedia.org/wiki/Students_for_Fair_Admissions_v._Harvard}{\textit{Students for Fair Admissions, Inc. v. President \& Fellows of Harvard College}, 600 U.S. 181 (2023)} (\href{https://www.everycrsreport.com/reports/RL30410.html}{CRS Report RL30410}; \href{https://scholar.smu.edu/context/smulr/article/4985/viewcontent/The_Post_Racial_Deception_of_the_Roberts_Court.pdf}{SMU Law Review}).
\item \textbf{Reserved Issue (Dicta/Assumption):} Powell expressly declined to decide whether Title VI implies a private right of action, assuming one only for purposes of the case (\href{https://supreme.justia.com/cases/federal/us/438/265/}{438 U.S. at 283--284}).
\end{itemize}

\section{Analysis of Opinions}
\subsection{Opinion Announcing the Judgment of the Court (Justice Powell)}
Justice Powell began by assuming without deciding that Bakke had a private right of action under Title VI and concluding from the 1964 legislative history that \S~601 of Title VI incorporates the constitutional standard of the Equal Protection Clause (\href{https://supreme.justia.com/cases/federal/us/438/265/}{438 U.S. at 283--287}).

Turning to the Fourteenth Amendment, Powell rejected the University's contention that ``benign'' racial preferences disadvantaging white applicants warrant less than strict scrutiny (\href{https://supreme.justia.com/cases/federal/us/438/265/}{438 U.S. at 287--299}). Because Fourteenth Amendment rights are ``personal rights'' guaranteed to every individual regardless of color, and because the United States is a ``Nation of minorities'' lacking any judicially manageable baseline for ranking which groups deserve preferred status, \textit{all} racial and ethnic classifications are inherently suspect and must satisfy strict scrutiny (\href{https://supreme.justia.com/cases/federal/us/438/265/}{438 U.S. at 289--297, 305}).

Powell then evaluated the four justifications advanced by UC Davis (\href{https://supreme.justia.com/cases/federal/us/438/265/}{438 U.S. at 306--315}):
\begin{enumerate}
\item \textit{Reducing the historic deficit of minorities in medicine:} Rejected as facially invalid ``discrimination for its own sake'' (\href{https://supreme.justia.com/cases/federal/us/438/265/}{438 U.S. at 307}).
\item \textit{Countering ``societal discrimination'':} Rejected as an amorphous concept absent judicial, legislative, or administrative findings of specific constitutional or statutory violations---findings a university faculty lacks competence to make (\href{https://supreme.justia.com/cases/federal/us/438/265/}{438 U.S. at 307--310}).
\item \textit{Improving health care in underserved communities:} Rejected for lack of empirical proof in the record that the racial quota actually promoted that goal (\href{https://supreme.justia.com/cases/federal/us/438/265/}{438 U.S. at 310--311}).
\item \textit{Attaining a diverse student body:} Accepted as a compelling state interest grounded in First Amendment academic freedom and the ``robust exchange of ideas'' (\href{https://supreme.justia.com/cases/federal/us/438/265/}{438 U.S. at 311--315}).
\end{enumerate}
However, Powell held that UC Davis's two-track 16-seat quota was not necessary to achieve genuine diversity because it insulated minority candidates from competition with all other applicants (\href{https://supreme.justia.com/cases/federal/us/438/265/}{438 U.S. at 315--320}). Pointing to the \textbf{Harvard College Admissions Program} (appended to his opinion), Powell concluded that a university may constitutionally treat race as one ``plus'' factor in individualized holistic review alongside geography, talents, and life experiences (\href{https://supreme.justia.com/cases/federal/us/438/265/}{438 U.S. at 316--318, 321--324}). Because UC Davis conceded it could not prove Bakke would have been rejected absent the unlawful quota, Bakke had to be admitted (\href{https://supreme.justia.com/cases/federal/us/438/265/}{438 U.S. at 320}).

\subsection{Joint Opinion Concurring in the Judgment in Part and Dissenting in Part (Justices Brennan, White, Marshall, and Blackmun)}
The Brennan Bloc agreed with Powell that Title VI's prohibition is coextensive with the Equal Protection Clause, emphasizing the 1964 legislative debates, HEW affirmative-action regulations (45 C.F.R. \S~80.3(b)(6)), subsequent congressional set-asides, and \textit{Lau v. Nichols}, 414 U.S. 563 (1974) (\href{https://supreme.justia.com/cases/federal/us/438/265/}{438 U.S. at 328--355}).

On the constitutional question, however, the four Justices argued that ``colorblindness'' is an aspiration rather than a description of reality and that benign remedial racial classifications should not trigger fatal strict scrutiny because white applicants are not a discrete and insular minority saddled with historical subordination (\href{https://supreme.justia.com/cases/federal/us/438/265/}{438 U.S. at 327, 355--358}). Borrowing from gender-discrimination precedent (\href{https://supreme.justia.com/cases/federal/us/429/190/}{\textit{Craig v. Boren}}), they proposed \textbf{intermediate scrutiny}: remedial racial classifications must serve important governmental objectives, be substantially related to those objectives, and avoid stigmatizing any group (\href{https://supreme.justia.com/cases/federal/us/438/265/}{438 U.S. at 359--362}). Applying that test, they would have upheld UC Davis's 16-seat program as a reasonable remedy for chronic minority underrepresentation in medicine resulting from past societal discrimination, and saw no constitutional difference between Davis's set-aside and Harvard's ``plus'' factor (\href{https://supreme.justia.com/cases/federal/us/438/265/}{438 U.S. at 362--379}).

\subsection{Separate Opinions by Justices White, Marshall, and Blackmun}
\begin{itemize}
\item \textbf{Justice White (438 U.S. at 379):} Wrote separately to contend that Title VI creates no implied private right of action, leaving enforcement exclusively to federal funding agencies under \S~602 (\href{https://tile.loc.gov/storage-services/service/ll/usrep/usrep438/usrep438265/usrep438265.pdf}{Library of Congress}; \href{https://caselaw.findlaw.com/court/us-supreme-court/438/265.html}{FindLaw}).
\item \textbf{Justice Marshall (438 U.S. at 387):} Delivered a sweeping historical account of slavery, \textit{Plessy v. Ferguson}, and Jim Crow, observing the deep irony that after centuries of class-based legal subjugation of Black Americans, the Court would invoke the Fourteenth Amendment to bar class-based remedies (\href{http://academic.brooklyn.cuny.edu/history/johnson/bakke.htm}{Brooklyn College Archive}; \href{https://aci.princeton.edu/document/65}{Princeton ACI}).
\item \textbf{Justice Blackmun (438 U.S. at 402):} Emphasized that universities have never admitted students solely on test scores---routinely favoring athletes, musicians, alumni children, and rural applicants---and famously argued that overcoming entrenched racism requires taking account of race (\href{https://liberty.lawbooks.cali.org/chapter/race-and-equal-protection-part-2/}{CALI}; \href{https://tcf.org/content/report/achieving-better-diversity/}{The Century Foundation}).
\end{itemize}

\subsection{Opinion Concurring in the Judgment in Part and Dissenting in Part (Justice Stevens, joined by Chief Justice Burger and Justices Stewart and Rehnquist)}
Invoking the canon of constitutional avoidance, the Stevens Bloc refused to reach the Fourteenth Amendment at all (\href{https://supreme.justia.com/cases/federal/us/438/265/}{438 U.S. at 408--421}; \href{https://en.wikipedia.org/wiki/Regents_of_the_University_of_California_v._Bakke}{Wikipedia}). Reading the California Supreme Court's judgment as limited to Bakke's own exclusion, Justice Stevens argued that the plain text of \S~601 of Title VI was ``crystal clear'' in barring any federally funded institution from excluding any person on the ground of race (\href{https://en.wikipedia.org/wiki/Regents_of_the_University_of_California_v._Bakke}{Wikipedia}; \href{https://supreme.justia.com/cases/federal/us/438/265/}{438 U.S. at 418}). Because UC Davis received federal funds and excluded Bakke from competing for 16 seats because of his race, the program violated Title VI, requiring affirmance of his admission and rendering any discussion of whether race could ever be used in admissions inappropriate (\href{https://en.wikipedia.org/wiki/Regents_of_the_University_of_California_v._Bakke}{Wikipedia}).

\section{Examples: Future Applications}
Evaluating future admissions and government-benefit policies under \textit{Bakke}'s own framework (while noting that \href{https://en.wikipedia.org/wiki/Students_for_Fair_Admissions_v._Harvard}{\textit{Students for Fair Admissions v. Harvard}, 600 U.S. 181 (2023)}, later effectively ended race-based college admissions):

\begin{enumerate}
\item \textbf{Same-Side (Unconstitutional under \textit{Bakke}): Automatic 20-Point Bonus for Minority Undergraduates.} A public university uses a 150-point admissions index and automatically awards 20 points to every applicant from an underrepresented racial minority group without individualized review of how each applicant contributes to diversity. \textit{Why:} Mirrors \href{https://supreme.justia.com/cases/federal/us/539/244/}{\textit{Gratz v. Bollinger}, 539 U.S. 244 (2003)}; like UC Davis's 16-seat quota, a mechanical, fixed racial bonus treats race as decisive rather than conducting the individualized, holistic comparison Powell required (\href{https://supreme.justia.com/cases/federal/us/438/265/}{438 U.S. at 317--320}).

\item \textbf{Same-Side (Unconstitutional under \textit{Bakke}): Separate Minority-Only Admission Track for Law School.} A public law school reserves 15\% of its entering class for a separate review subcommittee open exclusively to Black and Hispanic applicants, prohibiting white and Asian applicants from competing for those seats. \textit{Why:} Directly replicates UC Davis's two-track system; foreclosing non-preferred applicants from competing for a reserved percentage of seats violates individual Equal Protection rights under Powell's opinion and Title VI under Stevens's bloc (\href{https://supreme.justia.com/cases/federal/us/438/265/}{438 U.S. at 319--320}).

\item \textbf{Opposite-Side (Constitutional under \textit{Bakke}): Holistic, Individualized Review Treating Race as One ``Plus'' Factor.} A public law school evaluates every applicant in a single pool, weighing LSAT scores, GPAs, essays, work experience, overcoming adversity, and race as a flexible, non-mechanical ``plus'' factor to assemble a broadly diverse class, with no fixed numerical target. \textit{Why:} Directly tracks the Harvard College plan endorsed by Justice Powell (and later upheld in \href{https://en.wikipedia.org/wiki/Grutter_v._Bollinger}{\textit{Grutter v. Bollinger}, 539 U.S. 306 (2003)}, though later curtailed by \textit{SFFA} in 2023) (\href{https://supreme.justia.com/cases/federal/us/438/265/}{438 U.S. at 316--318}).

\item \textbf{Opposite-Side (Constitutional under \textit{Bakke}): Race-Neutral Socioeconomic and First-Generation Preference.} A public medical school grants a holistic ``plus'' factor to applicants of any race who grew up in low-income households, received Pell Grants, or are the first in their families to attend college. \textit{Why:} Does not classify on the basis of race at all, avoiding strict scrutiny and satisfying both Powell's individualized-consideration framework and the Stevens Bloc's colorblind reading of Title VI (\href{https://tcf.org/content/report/achieving-better-diversity/}{The Century Foundation}).

\item \textbf{Fence-Sitter (Close Case): Race-Neutral ``Top 10\%'' Plan Paired with Supplemental Holistic Review.} A state flagship university automatically admits the top 10\% of every public high school class (a race-neutral policy that captures racial diversity from geographically segregated schools) and uses race as a modest tiebreaker only for the remaining seats when race-neutral outreach has plateaued. \textit{Why:} This is a fence-sitter (echoing \href{https://supreme.justia.com/cases/federal/us/570/297/}{\textit{Fisher v. University of Texas}, 570 U.S. 297 (2013)}): under \textit{Bakke}, defenders argue the holistic component avoids quotas and fills gaps left by the percentage plan, while challengers argue that the success of the race-neutral 10\% plan proves that any additional racial classification is unnecessary and therefore fails narrow tailoring (\href{https://supreme.justia.com/cases/federal/us/570/297/}{570 U.S. at 312}).
\end{enumerate}

\section{Critique}
\textbf{Originalist and Textualist Perspectives:} Conservative, textualist, and originalist scholars sharply criticized both the Brennan Bloc and Justice Powell's compromise. In his famous 1979 essay \href{https://chicagounbound.uchicago.edu/cgi/viewcontent.cgi?article=11890&context=journal_articles}{\textit{The Disease as Cure}}, then-Professor Antonin Scalia attacked Justice Blackmun's claim that overcoming racism requires taking account of race, arguing that racial preferences perpetuate racial resentment, stigmatize minority achievers, and violate the textual principle of equal individual treatment (while endorsing race-neutral assistance for the poor) (\href{https://chicagounbound.uchicago.edu/cgi/viewcontent.cgi?article=11890&context=journal_articles}{Scalia, 1979 Wash. U. L.Q. 147, 153--156}). Textualists also echoed Justice Stevens's critique that Title VI's plain text (``No person \ldots shall, on the ground of race \ldots be excluded'') forbids \textit{any} race-based exclusion by federally funded universities, and argued that Powell's distinction between a rigid quota and a Harvard-style ``plus'' factor was a semantic veil that lacked grounding in constitutional text or original meaning (\href{https://www.journals.uchicago.edu/doi/full/10.1086/730860}{University of Chicago Press}; \href{https://www.everycrsreport.com/reports/RL30410.html}{CRS Report RL30410}). That anticlassification critique ultimately prevailed in \href{https://en.wikipedia.org/wiki/Students_for_Fair_Admissions_v._Harvard}{\textit{Students for Fair Admissions v. Harvard}, 600 U.S. 181 (2023)}.

\textbf{Progressive and Doctrinal Perspectives:} Progressive and antisubordination scholars have criticized \textit{Bakke} from the opposite direction---arguing that Justice Powell's opinion was actually a devastating loss for racial justice (\href{https://www.californialawreview.org/print/sffa-bakke}{California Law Review (2025)}; \href{https://www.californialawreview.org/print/antisubordination-eighth}{California Law Review (2024)}). By rejecting ``societal discrimination'' as a compelling interest and subjecting remedial affirmative action to the same strict scrutiny as Jim Crow segregation, Powell shifted Equal Protection from an \textit{antisubordination} guarantee (dismantling racial caste) to a symmetrical \textit{anticlassification} rule (\href{https://www.californialawreview.org/print/antisubordination-eighth}{California Law Review (2024)}). Critics argue that forcing universities to justify inclusion through ``classroom diversity'' rather than historical remediation treated minority students instrumentally---as educational assets for white peers---while leaving intact wealth- and legacy-based admissions preferences and laying the doctrinal foundation for \textit{SFFA} (\href{https://www.californialawreview.org/print/sffa-bakke}{California Law Review (2025)}; \href{https://academicworks.cuny.edu/cgi/viewcontent.cgi?article=1606&context=clr}{CUNY Law Review}).

\textbf{Underlying Competing Values:} \textit{Bakke} embodies the defining clash of modern Equal Protection law: \textit{formal individual equality and anticlassification} (protecting each person from being judged by race regardless of group history) versus \textit{substantive remedial equality and antisubordination} (permitting race-conscious measures to overcome entrenched, state-fostered racial disparities), mediated by Powell's pragmatic appeal to \textit{academic freedom and institutional pluralism} (\href{https://scholarship.law.wm.edu/cgi/viewcontent.cgi?article=3958&context=wmlr}{William \& Mary Law Review}).

\section{Key Quotations}
\begin{quote}
``The guarantee of equal protection cannot mean one thing when applied to one individual and something else when applied to a person of another color. If both are not accorded the same protection, then it is not equal.''
\end{quote}
\noindent--- Justice Powell, announcing the judgment of the Court, 438 U.S. at 289--290 (\href{https://supreme.justia.com/cases/federal/us/438/265/}{Justia}; \href{https://www.courtlistener.com/opinion/109930/regents-of-the-university-of-california-v-bakke/}{CourtListener}).

\begin{quote}
``Preferring members of any one group for no reason other than race or ethnic origin is discrimination for its own sake. This the Constitution forbids.''
\end{quote}
\noindent--- Justice Powell, 438 U.S. at 307 (\href{https://supreme.justia.com/cases/federal/us/438/265/}{Justia}; \href{https://constitutioncenter.org/the-constitution/supreme-court-case-library/university-of-california-v-bakke}{National Constitution Center}).

\begin{quote}
``In such an admissions program, race or ethnic background may be deemed a `plus' in a particular applicant's file, yet it does not insulate the individual from comparison with all other candidates for the available seats.''
\end{quote}
\noindent--- Justice Powell, 438 U.S. at 317 (\href{https://supreme.justia.com/cases/federal/us/438/265/}{Justia}; \href{https://www.law.cornell.edu/supremecourt/text/438/265}{Cornell LII}).

\begin{quote}
``In order to get beyond racism, we must first take account of race. There is no other way. And in order to treat some persons equally, we must treat them differently.''
\end{quote}
\noindent--- Justice Blackmun, concurring in the judgment in part and dissenting in part, 438 U.S. at 407 (\href{https://liberty.lawbooks.cali.org/chapter/race-and-equal-protection-part-2/}{CALI}; \href{https://chicagounbound.uchicago.edu/cgi/viewcontent.cgi?article=11890&context=journal_articles}{Chicago Unbound}).

\begin{quote}
``The meaning of the Title VI ban on exclusion is crystal clear: Race cannot be the basis of excluding anyone from a federally funded program.''
\end{quote}
\noindent--- Justice Stevens, joined by Chief Justice Burger and Justices Stewart and Rehnquist, concurring in the judgment in part and dissenting in part, 438 U.S. at 418 (\href{https://supreme.justia.com/cases/federal/us/438/265/}{Justia}; \href{https://en.wikipedia.org/wiki/Regents_of_the_University_of_California_v._Bakke}{Wikipedia}).

\end{document}
